My childhood was pretty much the stereotypical introvert childhood: a
passion for things like astronomy, barely speaking to anyone unless an
adult said it was mandatory, never really understanding social interaction.
what else can i say about my first 12 or 13 years of life? in reality, no
one really makes decisions below that age. some might argue that even
throughout our entire lives we're completely conditioned by our
environment, but at least during secondary school or high school we start
to make decisions: what degree to choose, what hobbies to pursue, who to
date, what music to listen to, what clothes to wear.

That's when i started thinking about what i would do. I started playing
the guitar, i started liking music outside of what my parents listened
to, and i started thinking seriously about life --- a bit philosophical for
a teenager, i guess. but when people ask me ``how did you do something?'',
the answer is usually: no friends and a lot of free time. And with that
two things you can start to do many things\ldots{} including getting crazy.

At the end of high school comes the moment where you say, ``ok, maybe i'll
choose something i'll do for the rest of my life''. i still think that if
i'd had a really good physics course, i might have become a physicist
instead of a chemist, but what's done is done --- i like chemistry, don't
get me wrong, they two aren't opposites. I also considered studying
philosophy, and sometimes i still wonder what my life would have been
like if i hadn't taken the science path.
